% ==============================================================================
% ssec:cs_posta --- Případová studie: KS Česká pošta 2025-2026
% Zdroje: KS_CeskaPosta_2025, CMKOS_ZpravaKV2025, CMKOS_PrinosKV2025,
%         zakon_zp_2006, zakon_kv_1991
% ==============================================================================

\Acf{KS} státního podniku \emph{Česká~pošta,~s.p.} na~léta 2025--2026 představuje jeden z~komplexních příkladů podnikového \aclgen{KV} v~českém prostředí a~současně zaměstnavatele zavázaného k~poskytování veřejných služeb.
Smlouva byla podepsána 30.~června 2025 v~Praze a~účinnost je sjednána od~1.~července 2025 do~30.~června 2026 s~automatickou prolongací do~31.~prosince 2027, pokud smluvní strany včas neuzavřou navazující \ac{KS}.
Z~hlediska motivací aktérů je~tato prolongace klíčová: oslabuje sice vyjednávací tlak v~následujícím kole, avšak státnímu zaměstnavateli zajišťuje provozní stabilitu a~zaměstnancům kontinuitu nároků, čímž odlišuje \acdat{KS} \acgen{ČP} od běžné roční účinnosti bez prolongace.\par

% Na~straně zaměstnanců vystupuje devět \aclpgen{OO}, jejichž společný postup koordinuje \acs{PKOV}~\acs{ČP} v~čele s~předsedkyní \acs{OS}~ZPTNS, sdruženého v~\acs{ČMKOS}.
% Ostaní \ac{OO} -- OSČMS, OS~DOSIA, poštovní sekce Federace strojvůdců, NOS~PP, OS~doručovatelů, OS~přepravy a~logistiky a~OS~administrativně-správních zaměstnanců \acgen{ČP},.
% Koordinační orgán není samostatným zákonným typem odborové organizace; prakticky však naplňuje požadavek \mbox{§~24} odst.~2 \aclgen{ZP}, aby při pluralitě odborových organizací u~jednoho zaměstnavatele jednaly všechny organizace společně a~ve~vzájemné shodě. Bez~takové shody by zaměstnavatel nemusel akceptovat samostatný návrh dílčí organizace jako společný návrh zaměstnanecké strany.
% Právě tato koordinovaná síla je~zdrojem vyjednávací pozice zaměstnanců: na~rozdíl od~maloobchodu, kde koalice především odvrací hrozbu blokády, dokáže poštovní koalice systematicky prosadit nadstandard napříč nárokovými položkami.\par

% Mzdová část se~opírá o~stupnici minimálních mezd v~jednotlivých tarifních stupních uvedenou v~Příloze~1, jejíž rozpětí pro standardní pracovní dobu \SI{40}{\hour\per\week} sahá od~přibližně \SI{20800}{\czk\per\month} v~prvním stupni do~\SI{42000}{\czk\per\month} ve~stupni dvanáctém.
% Nejnižší tarif platný od~poloviny roku 2025 leží pod úrovní zákonné minimální mzdy stanovené pro rok 2026 (\SI{22400}{\czk\per\month}), přičemž tato disparita nebyla v~mezisezónním kole adresována --- otevřená kompenzace je odložena na~vyjednávání nového znění \aclgen{KS} v~roce 2026.\par

Na straně zaměstnanců vystupuje devět odborových organizací. Podnikový koordinační odborový výbor, Česká pošta (\acs{PKOV}~\acs{ČP}), Podnikový výbor základních organizací Odborového svazu dopravy, České pošty~s.p. (PV~ZO~OSD~\acs{ČP}), Odborové hnutí zaměstnanců České pošty, s.p. (OHZ \acs{ČP}), SOS~--~21, nezávislé odborové sdružení pracovníků České pošty (SOS-21), Nezávislý odborový sněm České pošty (NOS \acs{ČP}), Základní odborová organizace \acs{OS} státních orgánů a~organizací, Krajské ředitelství policie Jihomoravského kraje (ZO~odborů PČR správa JM kraje), OSPEA (OSPEA), Federace výkonných zaměstnanců České republiky (FVZ), Svaz Odborářů Služeb a Dopravy (SOSaD). 
\acs{PKOV}~\acs{ČP} v čele s předsedkyní \acs{OS}~ZPTNS, sdruženého v~\acloc{ČMKOS} po změně \acgen{ZP}~\S~24 inicioval kolektivní vyjednávání mezi odborovými organizacemi pro přípravu společného návrhu \acgen{KS}, aby při pluralitě odborových organizací u~jednoho zaměstnavatele jednaly všechny organizace společně a~ve vzájemné shodě. Bez takové shody by zaměstnavatel mohl akceptovat samostatný návrh jedné či více organizací při dodržení podmínek \S~24~\acgen{ZP}. Právě koordinovaná síla je zdrojem vyjednávací pozice zaměstnanců, aby odborová  koalice mohla systematicky prosadit nadstandard napříč nárokovými položkami.\par

Mzdová část se opírá o~stupnici minimálních mezd v~jednotlivých tarifních stupních uvedenou v~Příloze~1, jejíž rozpětí pro standardní pracovní dobu 40~h/týd.~sahá od 20~800~Kč/měs. v prvním stupni (minimální mzda) do 42~000~Kč/měs. ve stupni dvanáctém.
Při změně minimální mzdy~\acs{ČR} se upraví minimální mzdy v jednotlivých tarifních stupních pro příslušný kalendářní rok tak, že budou mezi jednotlivými tarifními stupni dodrženy minimálně rozestupy uvedené v~tabulce „Výše minimálních mezd v~tarifních stupních“.\par

Naproti tomu mzdové příplatky překračují zákonné minimum napříč nárokovými položkami a~tvoří jádro protihodnoty pro zaměstnance.
Příplatek za~práci v~noci činí \SI{15}{\percent} \aclgen{PV} (minimálně \SI{26}{\czk\per\hour}) oproti zákonnému minimu \SI{10}{\percent} (\mbox{§~116} \aclgen{ZP}); příplatek za~práci v~sobotu a~neděli je sjednán na~\SI{12}{\percent} \aclgen{PV} (minimálně \SI{17}{\czk\per\hour}) oproti zákonným \SI{10}{\percent} (\mbox{§~118}); za~práci ve~ztíženém prostředí je sjednán plošný hodinový příplatek \SI{20}{\czk\per\hour} a~za~pracovní pohotovost \SI{15}{\percent} \aclgen{PV}.
Specifikem doručovatelské služby je smluvní zakotvení dělené směny s~odstupňovaným příplatkem \SIrange{26}{63}{\czk\per\hour} v~závislosti na délce přerušení, který český \aclgen{ZP} explicitně neupravuje a~který odpovídá typické poštovní organizaci práce: ranní výdej zásilek a~odpolední doručování.\par

Nejvýraznějším rysem smlouvy je~ovšem to, jak \aclgen{KV} přebírá rizika, která zákon i~veřejné sociální zabezpečení ponechávají na~jednotlivém pracovníkovi -- smlouva tak v~praxi naplňuje pojetí sociálního dialogu jako mechanismu sdílení rizik (kap.~\ref{sec:teorie}).
Prvním nástrojem je~pojištění odpovědnosti zaměstnanců za~škodu způsobenou zaměstnavateli, které zaměstnavatel sjednává a~plně hradí všem zaměstnancům. Tím zbavuje doručovatele a~přepážkové pracovníky regresního rizika, jež podle \aclgen{ZP} může u~schodku na~svěřených hodnotách dosáhnout násobků \aclgen{PV}.
Druhým nástrojem je~věrnostně odstupňované odstupné při skončení pracovního poměru z~organizačních důvodů (\S~52 písm.~a) až~c) \acs{ZP}) až do sedminásobku průměrného výdělku při trvání pracovního poměru nad dvacet pět let. Třetím nástrojem je odchodné při ztrátě zdravotní způsobilosti z~obecných důvodů (\mbox{§~52~e)} \aclgen{ZP}) až~do~čtyřnásobku \aclgen{PV} při trvání \aclgen{PP} nad dvacet pět let; zákon přitom v~tomto případě -- na~rozdíl od~ztráty způsobilosti následkem pracovního úrazu, žádné odchodné nepřiznává, takže smlouva pokrývá právě onen "dlouhý ocas" opotřebení ve~fyzicky náročných profesích, který sociální systém oceňuje jen omezeně a~který se~genderově dotýká zejména přepážkových pracovnic.
Tuto ochrannou logiku doplňuje dvoudenní zdravotní volno s~náhradou \SI{80}{\percent} \aclgen{PV}, třídenní volno s~rehabilitací při loupežném přepadení a~rozšířený doprovod zdravotně postiženého dítěte do~zařízení v~rozsahu šesti pracovních dnů ročně. Sjednán je příspěvek na produkty spoření na stáří ve výši 500~Kč, resp.~300~Kč po odpracování jednoho roku u~zaměstnavatele.\par
Dovolená je sjednána ve~výměře \SI{5}{\week}, tedy o~jeden týden nad zákonné minimum, stravování formou stravenkového paušálu \SI{100}{\czk} za~směnu (s~možností dvojnásobku při směně přesahující jedenáct hodin) a~náhrada za~použití vlastního jízdního kola doručovateli ve~výši \SI{4}{\czk\per\kilo\meter} představuje opět sektorově specifický nadstandard, který \aclgen{ZP} nezná a~který odpovídá ekonomické účelnosti jízdního kola v~podniku.
Benefitní rámec je doplněn sociálním fondem, z~něhož jsou financována další plnění zaměstnancům; u~státního podniku jde o~významný stabilizační prvek vedle samotných tarifních a~příplatkových ustanovení \aclgen{KS}.
V~oblasti vzdělávání jde Česká~pošta nad rámec maloobchodní praxe: zavazuje se zajišťovat prohlubování kvalifikace nezbytné k~výkonu sjednané práce a~při zvyšování kvalifikace poskytovat na~základě kvalifikační dohody pracovní úlevy a~hmotné zabezpečení; \acs{BOZP} má samostatnou část s~právem vyššího odborového orgánu kontrolovat její plnění, což spolu s~rehabilitací po~loupežném přepadení posouvá ochranu zdraví nad pouhou formální povinnost.
Souhrnně lze konstatovat, že vyjednávací bilance je~zde podstatně vyrovnanější než v~maloobchodě: silná koordinovaná koalice a~státní zaměstnavatel motivovaný spíše stabilitou a~sociálním smírem než minimalizací nákladů dosahují systematického nadstandardu nad~\aclgen{ZP}, čímž potvrzují empirické zjištění zprávy \acgen{ČMKOS} o~vyšší úrovni ochrany u~veřejnoprávních zaměstnavatelů s~koordinovanou odborovou stranou.~\cite{KS_CeskaPosta_2025}~\cite{CMKOS_ZpravaKV2025}~\cite{CMKOS_PrinosKV2025}\par
